\section{Raw data}\label{sec:rawdata}

\subsection{Constants of the apparatus}

\begin{table}[htbp]
\centering\footnotesize
\caption{Constants of the apparatus. Values in red are still to be measured (Section~\ref{sec:method}).}
\label{tab:constants}
\setlength{\tabcolsep}{4pt}
\begin{tabular}{@{}p{4.3cm}cp{4.6cm}p{3.6cm}@{}}
\toprule
Quantity & Symbol & Value & Source \\
\midrule
Free length of spring & $L$ & \TBM{\ldots}$\pm0.5\un{mm}$ & rule \\
Width, thickness of strip & $b$, $h$ & \TBM{\ldots}, \TBM{\ldots}$\un{mm}$ ($\pm0.02$) & calliper (mean of 5) \\
Mass per unit length & $\mu$ & \TBM{\ldots}$\un{g\,m^{-1}}$ & balance, 300 mm of strip \\
Magnet diameter, thickness & $D$, $t$ & \TBM{\ldots}, \TBM{\ldots}$\un{mm}$ & calliper \\
Magnet mass (with holder) & $M$ & \TBM{\ldots}$\pm0.01\un{g}$ & balance \\
Static stiffness, spring 1 / 2 & $k$ & \TBM{\ldots} / \TBM{\ldots}$\un{N\,m^{-1}}$ & load--deflection graph \\
Uncoupled frequency, spring 1 / 2 & $f_a$ / $f_b$ & \TBM{\ldots} / \TBM{\ldots}$\un{Hz}$ ($\pm0.003$) & FFT, 60 s, 3 trials \\
Effective mass & $m=k/(2\pi f_a)^2$ & \TBM{\ldots}$\un{g}$ & derived \\
Magnetic moment & $\mum$ & \TBM{\ldots}$\un{A\,m^2}$ & balance method / datasheet \\
Predicted constant & $C=3\muz\mum^2/(\pi^3m)$ & \TBM{\ldots}$\un{m^5\,s^{-2}}$ & derived \\
Decay time of oscillation & $\tau$ & $\approx25\un{s}$ & envelope of $x_1(t)$ \\
Frame rate, record length & $f_s$, $T$ & $240\un{fps}$, $30\un{s}$ ($N=\nNsamples$) & camera metadata \\
\bottomrule
\end{tabular}
\end{table}

\begin{table}[htbp]
\centering\small
\caption{Static load--deflection data for spring 1 (spring 2 in Appendix~\ref{app:rawdata}). \TBM{to be measured}}
\label{tab:static}
\begin{tabular}{@{}cccccc@{}}
\toprule
load / g & $F$ / mN & $x$ (loading) / mm & $x$ (unloading) / mm & mean $x$ / mm & $\Delta x$ / mm\\
\midrule
1 & 9.8 & \TBM{} & \TBM{} & \TBM{} & 0.5\\
2 & 19.6 & & & & \\
$\vdots$ & & & & & \\
10 & 98.1 & & & & \\
\bottomrule
\end{tabular}
\end{table}

\subsection{Mode frequencies}

Table~\ref{tab:rawfreq} lists the two peak frequencies read from the spectrum of each trial. All spectra were computed from $T=30\un{s}$ records sampled at $f_s=240\un{Hz}$ ($N=\nNsamples$ samples), detrended, with a rectangular window, zero-padded to $2^{20}$ points and with parabolic peak interpolation; the bin width is $1/T=\nbinT\un{mHz}$. The uncertainty of each mean is taken as the half-range of the three trials (three trials are too few for a meaningful standard deviation) or the resolution-limited value $\pm5\un{mHz}$ (Section~\ref{sec:fftres}), whichever is larger. The last column states whether the two peaks were resolved, using the criterion $f_2-f_1\geq2/T=67\un{mHz}$ for a clean separation with a rectangular window; the cases $35$ and $40\un{mm}$ are marked as \emph{marginal}, and at $45$, $50$ and $60\un{mm}$ only a single peak was present. The unresolved rows are a result in themselves (they mark where the beat period exceeds the record length) rather than missing data.

\begin{table}[htbp]
\centering\footnotesize
\caption{Peak frequencies read from the FFT of $x_1(t)$ for each trial ($A=10\un{mm}$, $T=30\un{s}$, $240\un{fps}$, rectangular window). The trial columns are to be filled from the raw exports; the means are the values used in the analysis.}
\label{tab:rawfreq}
\setlength{\tabcolsep}{3.5pt}
\begin{tabular}{@{}c ccc c ccc c c@{}}
\toprule
$d\pm0.5$ & \multicolumn{3}{c}{$f_1$ / Hz} & mean $f_1$ & \multicolumn{3}{c}{$f_2$ / Hz} & mean $f_2$ & resolved?\\
/ mm & tr.\,1 & tr.\,2 & tr.\,3 & $\pm0.005$ & tr.\,1 & tr.\,2 & tr.\,3 & $\pm0.005$ & ($f_2-f_1\geq2/T$)\\
\midrule
21 & \TBM{} & \TBM{} & \TBM{} & 0.823 & \TBM{} & \TBM{} & \TBM{} & 1.063 & yes\\
23 & & & & 0.808 & & & & 1.005 & yes\\
25 & & & & 0.812 & & & & 0.953 & yes\\
28 & & & & 0.810 & & & & 0.907 & yes\\
31 & & & & 0.807 & & & & 0.881 & yes\\
35 & & & & 0.806 & & & & 0.845 & marginal (1.2 bins)\\
40 & & & & 0.802 & & & & 0.811 & no (0.3 bins)\\
45 & & & & \multicolumn{5}{c}{single peak at $\approx0.80\un{Hz}$} & no\\
50 & & & & \multicolumn{5}{c}{single peak at $\approx0.80\un{Hz}$} & no\\
60 & & & & \multicolumn{5}{c}{single peak at $\approx0.80\un{Hz}$} & no\\
\bottomrule
\end{tabular}
\end{table}

\paragraph{Qualitative observations.} At $21$--$25\un{mm}$ the beats were obvious to the eye, with the energy passing completely from one spring to the other and back within a few seconds. At $35\un{mm}$ one beat took roughly $25\un{s}$, comparable to the time for the amplitude to decay by a factor $e$, and by $45\un{mm}$ the second spring barely moved during the whole record. A small amount of motion perpendicular to the line joining the magnets was visible when the release was not perfectly horizontal; those trials were repeated. The displacement traces also showed that the oscillation frequency of the anti-phase motion was slightly higher during the first few seconds (largest amplitude) than later, an observation that becomes important in Section~\ref{sec:nonlinear}.
